\uuid{Gn1Y}
\titre{Points critiques et extremums de $f(x,y)=x^2y^3+(y-1)^2$}
\niveau{L2}
\module{Analyse}
\chapitre{Fonctions de plusieurs variables}
\sousChapitre{Optimisation sans contrainte}
\theme{Points critiques, hessienne, minimum local, minimum global}
\auteur{Maxime Nguyen}
\datecreate{2026-03-24}
\organisation{amscc}
\difficulte{3}

\contenu{
\texte{
  On considère la fonction $f : \mathbb{R}^2 \to \mathbb{R}$ définie par $f(x,y) = x^2y^3 + (y-1)^2$.
}

   \question{Déterminer les points critiques de $f$. Admet-elle des extremums locaux ? Les extremums locaux trouvés sont-ils globaux ?}
    \reponse{
      On calcule :
      \[
        \frac{\partial f}{\partial x} = 2xy^3, \qquad \frac{\partial f}{\partial y} = 3x^2y^2 + 2(y-1).
      \]
      \begin{itemize}
        \item $\dfrac{\partial f}{\partial x} = 0 \Rightarrow x = 0$ ou $y = 0$.
        \item Cas $y = 0$ : $\dfrac{\partial f}{\partial y} = -2 \neq 0$ $\Rightarrow$ pas de point critique.
        \item Cas $x = 0$ : $\dfrac{\partial f}{\partial y} = 2(y-1) = 0 \Rightarrow y = 1$.
      \end{itemize}
      L'unique point critique est $(0,1)$.

      Les dérivées secondes en $(0,1)$ :
      \[
        \frac{\partial^2 f}{\partial x^2}(0,1) = 2, \quad \frac{\partial^2 f}{\partial y^2}(0,1) = 2, \quad \frac{\partial^2 f}{\partial x\partial y}(0,1) = 0.
      \]
      \[
        \operatorname{Hess}_f(0,1) = \begin{pmatrix} 2 & 0 \\ 0 & 2 \end{pmatrix}, \quad \Delta_f(0,1) = 4 > 0, \quad \frac{\partial^2 f}{\partial x^2} > 0.
      \]
      Donc $f$ admet un minimum local en $(0,1)$. Cependant, $f(1,y) \to -\infty$ quand $y \to -\infty$, donc ce minimum local n'est pas global.
    }

}
